\chapter{Thảo luận, phòng vệ và kết luận}
\label{chap:conclusion}

\section{Ý nghĩa về quyền riêng tư}

Thực nghiệm cho thấy một thông điệp rõ ràng: ẩn danh hóa không đồng nghĩa với riêng tư. Việc gỡ bỏ tên và số định danh khỏi một tập micro-data thưa gần như không bảo vệ được người dùng, vì chính tổ hợp hành vi của họ đã là một ``dấu vân tay'' duy nhất. Hệ quả còn nghiêm trọng hơn với đặc tính forward privacy: một khi một mẩu dữ liệu bị liên kết với danh tính thật, mọi liên kết về sau (kể cả với các định danh ảo mới) đều trở nên dễ dàng hơn. Trong bối cảnh thiết bị và ứng dụng di động liên tục thu thập dữ liệu vị trí, lịch sử duyệt web, danh sách ứng dụng đã cài - vốn cũng thưa và nhiều chiều - rủi ro giải ẩn danh là hoàn toàn hiện hữu.

\section{Các biện pháp phòng vệ}

\subsection{Vì sao k-anonymity không đủ}
Như đã phân tích ở Chương~\ref{chap:background}, k-anonymity giả định tách được quasi-identifier khỏi thuộc tính nhạy cảm - điều bất khả thi với dữ liệu hành vi, nơi mọi thuộc tính đều có thể là điểm liên kết. Trên dữ liệu nhiều chiều, để đạt $k$-ẩn danh phải khái quát hóa tới mức dữ liệu mất giá trị sử dụng.

\subsection{Riêng tư vi phân (Differential Privacy)}
Hướng phòng vệ vững chắc nhất hiện nay là riêng tư vi phân~\cite{dwork}: thay vì công bố micro-data, hệ thống chỉ trả lời các truy vấn tổng hợp và cố ý thêm nhiễu được hiệu chỉnh toán học, sao cho sự hiện diện hay vắng mặt của \textit{một} cá nhân gần như không ảnh hưởng tới kết quả. Cách này cung cấp bảo đảm riêng tư có thể chứng minh được, đánh đổi bằng một phần độ chính xác.

\subsection{Các biện pháp thực tiễn khác}
\renewcommand{\labelitemi}{$-$}
\begin{itemize}
	\item \textbf{Không công bố micro-data ở mức cá nhân}; chỉ phát hành thống kê tổng hợp đã thêm nhiễu.
	\item \textbf{Thêm nhiễu mạnh và lấy mẫu nhỏ:} tuy nhiên bài báo chứng minh nhiễu nhẹ là vô tác dụng - phải đủ mạnh tới mức ảnh hưởng giá trị sử dụng mới có hiệu quả.
	\item \textbf{Giảm chiều và gộp nhóm thuộc tính:} loại bỏ hoặc gộp các thuộc tính hiếm (vốn mang giá trị định danh cao nhất) trước khi công bố.
	\item \textbf{Kiểm soát truy cập và ràng buộc pháp lý:} hợp đồng sử dụng dữ liệu, cấm tái định danh, tối thiểu hóa dữ liệu thu thập ngay từ đầu.
\end{itemize}

\section{Hạn chế của nghiên cứu}

\renewcommand{\labelitemi}{$-$}
\begin{itemize}
	\item Tập MovieLens 100k nhỏ hơn nhiều so với Netflix (943 so với 500.000 người), nên đám đông ứng viên ít hơn và một số con số tuyệt đối lạc quan hơn thực tế; tuy nhiên các \textit{xu hướng} được tái hiện trung thực.
	\item Các tham số $\rho_0, d_0, \phi$ lấy theo bài báo, chưa tinh chỉnh riêng cho MovieLens.
	\item Hệ thống mô phỏng kịch bản kẻ tấn công đã biết bản ghi nạn nhân nằm trong tập công bố; chưa cài đặt đầy đủ biến thể phát hiện ``nạn nhân không có trong mẫu''.
\end{itemize}

\section{Kết luận}

Báo cáo đã nghiên cứu, mô phỏng và kiểm chứng bằng thực nghiệm thật lớp tấn công giải ẩn danh trên dữ liệu thưa của Narayanan-Shmatikov. Chúng em đã: (1) trình bày lại nền tảng lý thuyết và thuật toán Scoreboard-RH; (2) xây dựng, sửa lỗi và nâng cấp một ứng dụng web hoàn chỉnh - đặc biệt là sửa hàm cho điểm cho đúng công thức gốc, bổ sung chế độ không-ngày, phơi bày trọng số độ hiếm và mô-đun thực nghiệm; (3) tái hiện thành công các kết quả chính của bài báo trên MovieLens 100k, với 8 phim cho tỉ lệ định danh đúng tới 98\%, và chứng minh tính bền vững trước nhiễu cũng như sức mạnh đặc biệt của các thuộc tính hiếm.

Kết quả khẳng định một bài học nền tảng của an toàn dữ liệu hiện đại: trong thế giới dữ liệu thưa và nhiều chiều, riêng tư phải được bảo đảm bằng các cơ chế có cơ sở toán học như riêng tư vi phân, chứ không thể trông cậy vào việc ``gỡ bỏ thông tin định danh''.
